\documentclass[10pt,a4paper]{amsart}
\usepackage[margin=19mm]{geometry}
\usepackage{amsmath,amssymb,mathtools}
\usepackage[scr=rsfs,cal=euler]{mathalfa}
\usepackage{xfrac}
\usepackage[unicode,hidelinks,bookmarks=false]{hyperref}
\newcommand{\ie}{i.e.\ }
\newcommand{\cf}{cf.\ }
\newcommand{\resp}{respectively\ }
\newcommand{\Subsection}{\subsection}
\providecommand{\nobreakdash}{\nobreak\hbox{-}}
\setlength{\emergencystretch}{2em}
\multlinegap0pt
\usepackage{bm}
\usepackage{bbm}
\usepackage{dsfont}
\usepackage{xspace}
\usepackage{cancel}
\usepackage{mathtools}
\usepackage{cleveref}
\usepackage{amsthm}
\makeatletter
\@ifundefined{theorem}{\newtheorem{theorem}{Theorem}}{}
\@ifundefined{claim}{\newtheorem{claim}[theorem]{Claim}}{}
\@ifundefined{corollary}{\newtheorem{corollary}[theorem]{Corollary}}{}
\@ifundefined{lemma}{\newtheorem{lemma}[theorem]{Lemma}}{}
\makeatother
\title[Linear Tur\'an Numbers of Uniform Hypertrees]{Linear Tur\'an Numbers of Uniform Hypertrees}
\author{Rajat Adak, Pragya Verma}
\date{Source: arXiv:2607.16854v2 (CC BY 4.0)}
\begin{document}
\maketitle

\noindent\emph{Source and licence.} Linear Tur\'an Numbers of Uniform Hypertrees. Version arXiv:2607.16854v2 (CC BY 4.0), distributed under the Creative Commons Attribution 4.0 International licence, as recorded in the arXiv licence metadata for this exact version and shown on the abstract page https://arxiv.org/abs/2607.16854v2. Full rights notice: licenses/arxiv-2607.16854-CC-BY-4.0.txt.\par\smallskip

\noindent\textit{Editorial scope.} Counted unit: the Lemma whose source label is \texttt{None} in the article (source0.tex lines 1474--1481, source label \texttt{None}), delivered together with the article's own complete original proof of it (source0.tex lines 1483--1623, one \texttt{proof} environment of 141 lines). No mathematics is written, repaired or re-derived: statement and proof are the article's own text line for line. Required context retained uncounted: lem:path-line-graph (lines 1356--1360); joincor (lines 1447--1455). Every definition, display and label of those lines is retained unchanged. Omitted from this package and not redistributed: 1--1355, 1361--1446, 1456--1473, 1482--1482, 1624--2410. Nothing inside a retained excerpt is omitted or altered: every retained body is delivered exactly as the article's lines are written, so no omission receipt of any kind accompanies this package. Citations: the retained excerpts carry no citation command at all, so no bibliography entry of the article is transcribed and no reference is redistributed. Reference adaptation: none was needed and none was made; every reference inside the delivered unit resolves inside it, so no unresolved reference or citation is produced. Printed slips kept literal: no printed word, symbol, hypothesis or number of the counted statement or of its proof was altered. Layout and class: the article's own class is \texttt{article} with packages including fontenc, inputenc, geometry, amsmath, amssymb, amsfonts, amsthm, graphicx, tikz, float, capt-of, comment, natbib, hyperref; the wrapper is the lane's pinned \texttt{amsart} editorial wrapper at 10pt on A4, which restates the theorem-like environments the retained text needs and adds the article's own macro definitions that the retained text needs (none), with extra wrapper packages (bm, bbm, dsfont, xspace, cancel, mathtools, cleveref). The delivered numbering is the unit's own. Assets: no retained span contains a figure environment, a graphics-inclusion command, a PDF-page inclusion, a table or a TikZ picture; no image file is copied into this package, the package declares no asset dependency and no image file is redistributed with it. Licence: version arXiv:2607.16854v2 of this article is distributed under the Creative Commons Attribution 4.0 International licence, as recorded in the arXiv licence metadata for this exact version and shown on the abstract page https://arxiv.org/abs/2607.16854v2. A full rights notice is delivered at licenses/arxiv-2607.16854-CC-BY-4.0.txt. Characters: every retained excerpt is pure ASCII, so the delivered engine, XeLaTeX with its default Latin Modern text font, reports no missing character.
\begin{lemma}\label[lemma]{lem:path-line-graph}
Let $H$ be a linear $r$-uniform hypergraph. Four distinct hyperedges
$f_1,f_2,f_3,f_4$ form a copy of $P_4^r$, in this order, if and only if
the corresponding vertices induce a path on four vertices in $L(H)$.
\end{lemma}

\begin{corollary}
\label[corollary]{joincor}
Let $G$ be a connected graph with at least two vertices and containing
no induced path on four vertices, that is $G$ is a connected cograph. Then there exists a nontrivial partition
    $V(G)=A\mathbin{\dot\cup}B$ such that every vertex of $A$ is adjacent to every vertex of $B$.
Equivalently,
    $$G=G[A]\vee G[B],$$
where $\vee$ denotes the graph join.
\end{corollary}

\begin{lemma}\label[lemma]{joinineq}
Let $H$ be a connected linear $r$-uniform $P_4^r$-free hypergraph. Let
$
m=|E(H)|$ and $k=\Delta(H).$
Then $m\leq rk.$
Moreover, if $m=rk$, then every nonisolated vertex of $H$ has degree
$k$.
\end{lemma}

\begin{proof}
If $m=1$, then $k=1$ and the inequality follows from $r\geq 2$. Also,
equality is impossible in this case. Thus, assume that $m\geq 2$.

By \Cref{lem:path-line-graph}, $L(H)$ contains no induced path on four
vertices. Since $H$ is connected, $L(H)$ is connected. By
\Cref{joincor}, there exists a partition
$$
E(H)=\mathcal{A}\sqcup\mathcal{B},
\qquad
|\mathcal{A}|=s,
\qquad
|\mathcal{B}|=t,
$$
where $s,t>0$, such that every hyperedge of $\mathcal{A}$ intersects every
member of $\mathcal{B}$. For each $x\in V(H)$, define
$$
a_x=|\{A\in\mathcal{A}:x\in A\}|,
\qquad
b_x=|\{B\in\mathcal{B}:x\in B\}|,
\qquad
d_H(x)=a_x+b_x.
$$
Counting incidences with the two edge classes gives
$$
\sum_{x\in V(H)}a_x=rs
\qquad\text{and}\qquad
\sum_{x\in V(H)}b_x=rt.
$$
Every pair of hyperedges $(A,B)\in\mathcal{A}\times\mathcal{B}$ has exactly one
common vertex: it has at least one common vertex by the choice of the
partition, and at most one by linearity. Hence,
$$
\sum_{x\in V(H)}a_xb_x=st.
$$

We now compare the number of cross-pairs between $\mathcal{A}$ and
$\mathcal{B}$ with the maximum degree of $H$.
Since, $\sum_{x\in V(H)}a_xb_x=st$, $a_xb_x$ can be viewed as the number of pairs
$(A,B)\in\mathcal{A}\times\mathcal{B}$ whose unique intersection occurs
at the vertex $x$.

Recall that $\Delta(H) = k$. Our aim is to bound $a_xb_x$ locally using the degree
    $d_H(x)=a_x+b_x\leq k.$
For this purpose, put
    $p=\dfrac{s}{m}$ and $q=\dfrac{t}{m}.$
Thus $p+q=1$. 

The choice of these two quantities is motivated by the fact
that the total $\mathcal{A}$ and $\mathcal{B}$ incidences are
    $\sum_x a_x=rs$ and $\sum_x b_x=rt,$
and we want to obtain the term $rst/m$ after summing a suitable linear
combination of $a_x$ and $b_x$.

\begin{claim}\label[claim]{cl6.5}
    For a vertex $x$ with $d_H(x) >0$,  $\dfrac{a_xb_x}{d_H(x)}
    \leq
    q^2a_x+p^2b_x.$
\end{claim}
   \begin{proof} Recall that $p+q =1$, therefore $p^2+q^2 -1 = -2pq$. Thus
\begin{align*}
q^2a_x+p^2b_x-\frac{a_xb_x}{d_H(x)}
&=
q^2a_x+p^2b_x-\frac{a_xb_x}{a_x+b_x}
=
\frac{
q^2a_x(a_x+b_x)
+
p^2b_x(a_x+b_x)
-
a_xb_x
}{d_H(x)}\\
&=
\frac{
q^2a_x^2+p^2b_x^2+
(q^2+p^2-1)a_xb_x
}{d_H(x)}\\
&=
\frac{q^2a_x^2-2pqa_xb_x+p^2b_x^2}{d_H(x)}
=
\frac{(qa_x-pb_x)^2}{d_H(x)}
\geq 0.
\end{align*}\end{proof}
Now, since $d_H(x)\leq k$, we also have
    $\dfrac{a_xb_x}{k}
    \leq
    \dfrac{a_xb_x}{d_H(x)}.$
Thus from \Cref{cl6.5} we get
$$
    \frac{a_xb_x}{k}
    \leq
    \frac{a_xb_x}{d_H(x)}
    \leq
    q^2a_x+p^2b_x.
$$
Thus, summing over all vertices with $d_H(x)>0$ we get,
$$
\frac{1}{k}\sum_{x:d_H(x)>0}a_xb_x
\leq
q^2\sum_{x:d_H(x)>0}a_x
+
p^2\sum_{x:d_H(x)>0}b_x=
q^2rs+p^2rt.
$$
Using $\sum_xa_xb_x=st,$ we obtain
    $\dfrac{st}{k}
    \leq
    q^2rs+p^2rt.$
We now simplify the right-hand side. Since
    $p=\dfrac{s}{m},$ and $q=\dfrac{t}{m},$
and $s+t=m$, we have
$$
q^2rs+p^2rt
=
r\left(
\frac{t^2s}{m^2}
+
\frac{s^2t}{m^2}
\right)=
\frac{rst(s+t)}{m^2}=
\frac{rst}{m}.
$$
Therefore, $\dfrac{st}{k}\leq\dfrac{rst}{m}.$
Since $s,t>0$, we may cancel $st$, obtaining
$\dfrac{1}{k}\leq\dfrac{r}{m}.$ Therefore,
    $m\leq rk.$
    
\vspace{2mm} 

For the equality case, observe that, for every non-isolated vertex $x$,
$$
q^{2}a_x+p^{2}b_x-\frac{a_xb_x}{k}
=
\frac{(qa_x-pb_x)^{2}}{d_H(x)}
+a_xb_x\left(\frac{1}{d_H(x)}-\frac{1}{k}\right).
$$
Both terms on the right are nonnegative. If $m=rk$, equality forces both
terms to vanish. Thus $qa_x=pb_x$. Since $p,q>0$ and
$a_x+b_x=d_H(x)>0$, we have $a_xb_x>0$. Consequently,
$a_xb_x\left(\dfrac{1}{d_H(x)}-\dfrac{1}{k}\right)=0$ implies $d_H(x)=k$. Hence $H$ is $k$-regular.
\end{proof}

\end{document}
